\section{Overview}
The default implementation of Kyber in \texttt{pqm4} is designed to be resitant to fault injection attacks. To create a controlled environment for testing and developing the attack, we introduced specific modifications to the codebase. These changes were made primarily in \texttt{crypto\_kem/ml-kem-768/m4fspeed/kem.c} and \texttt{crypto\_kem/ml-kem-768/m4fspeed/verify.c}.

\section{Modifications in `kem.c`}
The core vulnerability was introduced by modifying the decapsulation logic to use a "default-to-true" pattern. In a secure implementation, the shared secret should only be copied if the validation succeeds. We modified this to first copy the true shared secret, and then overwrite it with a rejection key if validation fails.

This creates a race condition: if we can skip the overwrite instruction, the device will return the valid shared secret even for an invalid ciphertext.

\subsection{Code Snippet}
The critical section in \texttt{crypto\_kem\_dec} function was modified as follows:

\begin{lstlisting}[language=C]
    /* Vulnerable Pattern: Default to True Key, overwrite if fail */
    memcpy(ss, kr, KYBER_SYMBYTES);

    /* Overwrite with rejection key if fail is true */
    hal_trigger_toggle();
    cmov(ss, rejection_key, KYBER_SYMBYTES, fail);
    hal_trigger_toggle();
\end{lstlisting}

We also added \texttt{hal\_trigger\_toggle()} calls around the \texttt{cmov} function. These triggers allow us to synchronize the ChipWhisperer glitch with the exact execution window of the target instruction.

\section{Modifications in `verify.c` (Target Instruction)}
The \texttt{cmov} function handles the conditional move. We targeted the specific instruction responsible for negating the condition flag \texttt{b}.

\begin{lstlisting}[language=C]
void cmov(unsigned char *r, const unsigned char *x, size_t len, unsigned char b) {
    size_t i;

    b = -b;  // <-- This instruction is the GLITCH TARGET
    for (i = 0; i < len; i++) {
        r[i] ^= b & (x[i] ^ r[i]);
    }
}
\end{lstlisting}


\section{Build System Modifications}
To facilitate the creation of the specific test firmware, we modified the build system by adding a new configuration file \texttt{mk/tests.mk}. This file defines the \texttt{test-glitch} target and lists all necessary source files, including the hardware abstraction layer and the Kyber implementation files.

\begin{lstlisting}[language=make]
# Custom test for glitch attack
TEST_GLITCH_C_SOURCES = \
    test_glitch.c \
    crypto_kem/ml-kem-768/m4fspeed/cbd.c \
    ...
    common/hal-opencm3.c \
    common/randombytes.c

elf/test_glitch.elf: $(TEST_GLITCH_OBJS) $(LINKDEPS) $(CONFIG)
\end{lstlisting}

This allows us to compile the attack firmware using the standard \texttt{make} command:
\texttt{make PLATFORM=cw308t-stm32f3 test-glitch}.

\section{Deterministic Randomness}
For a controlled fault injection experiment, it is essential that the device behaves deterministically. Standard Kyber implementations use a hardware random number generator (TRNG) or a seeded PRNG to generate ephemeral keys and noise.

We modified \texttt{common/randombytes.c} to use a \textbf{fixed seed} derived from the digits of $\pi$. This ensures that every execution of the key generation and encapsulation process produces the exact same "random" values, allowing us to:
\begin{enumerate}
    \item Knowing the "Gold" (correct) keys in advance.
    \item Reliably reproduce the exact same state for every glitch attempt.
    \item Differentiate between a successful glitch (returning the Gold key) and a device reset or random failure.
\end{enumerate}

\begin{lstlisting}[language=C]
// Fixed seed for deterministic execution (Pi digits)
static uint32_t seed[32] = {2, 1, 4, 1, 5, 9, ...};
\end{lstlisting}

